\documentclass[12pt]{exam}

\usepackage[utf8]{inputenc}
\usepackage[T1]{fontenc}
\usepackage{lmodern}
\usepackage[margin=1in]{geometry}
\usepackage{amsmath,amssymb}
\usepackage{array}
\usepackage{enumitem}
\usepackage{fancyvrb}

\renewcommand{\thepartno}{\alph{partno}}
\renewcommand{\partlabel}{\thepartno.}
\renewcommand{\questionlabel}{\thequestion.}

\pagestyle{headandfoot}
\header{Name \rule[-2pt]{6.5cm}{0.4pt}}{}{SID \rule[-2pt]{6cm}{0.4pt}}
\footer{Spring 2026}{\thepage}{CS/EE 147 Practice Exam 2}
\headrule

\newcommand{\Ans}{\underline{Answer:}\ }
\newcommand{\smallbox}{\framebox[5cm]{\rule{0pt}{1.1cm}}}
\newcommand{\tfbox}{\framebox[1.7cm]{\rule{0pt}{1.7cm}}}
\newcommand{\widebox}[1][9.5cm]{\framebox[#1]{\rule{0pt}{1.1cm}}}
\newcommand{\paragraphbox}[1]{\framebox[\linewidth]{\rule{0pt}{#1}}}
\newcommand{\mc}{$\bigcirc$\ }

\setlength{\parindent}{0pt}

\begin{document}

% =====================================================================
%  TITLE PAGE
% =====================================================================
\begin{center}
  {\Large\bfseries CS/EE 147}\\[2pt]
  {\Large\bfseries GPU Architecture and Parallel Programming}\\[18pt]
  {\large\bfseries Spring 2026}\\[14pt]
  {\large\bfseries Practice Exam 2}\\[14pt]
  {\large\bfseries Time Allowed: 80 Minutes}
\end{center}

\vspace{10pt}

\begin{center}
\renewcommand{\arraystretch}{1.45}
\begin{tabular}{|>{\centering\arraybackslash}m{4cm}|>{\centering\arraybackslash}m{4cm}|>{\centering\arraybackslash}m{4cm}|}
\hline
\textbf{Question} & \textbf{Points Possible} & \textbf{Points Scored} \\ \hline
1            & 4  & \\ \hline
2            & 4  & \\ \hline
3            & 4  & \\ \hline
4            & 13 & \\ \hline
5            & 12 & \\ \hline
6--9 (MC)    & 12 & \\ \hline
10--15 (T/F) & 12 & \\ \hline
16           & 6  & \\ \hline
17           & 12 & \\ \hline
18           & 4  & \\ \hline
Total        & 83 & \\ \hline
\end{tabular}
\end{center}

\vspace{14pt}

\begin{itemize}[leftmargin=2.2em,itemsep=4pt]
  \item Closed book, closed internet
  \item Allowed 1 page letter-size cheatsheet, front and back
  \item Calculators allowed, but likely not necessary
  \item Covers Lectures 8--14: Unified Memory, Matrix Multiply, Memory Coalescing, Histogram, Multi-GPU, GPU Sharing (slides 1--25), Dynamic Parallelism
\end{itemize}

\newpage

\begin{questions}

% ---------------- Q1 ----------------
\question \textbf{[4 points] Tiled SGEMM Launch}\\
You launch a tiled matrix-multiply kernel on a $1024 \times 1024$ output matrix with
\texttt{TILE\_WIDTH = 32}.
\begin{parts}
  \part {}[2 points] How many warps are in each block?\\[8pt]
        \Ans \smallbox
  \part {}[2 points] How many blocks are in the grid?\\[8pt]
        \Ans \smallbox
\end{parts}

\vspace{6pt}
% ---------------- Q2 ----------------
\question \textbf{[4 points]} Below is the loop body of an interleaved-partition histogram
kernel. What should \texttt{stride} be initialized to so that all threads in the grid jointly
cover the input with a coalesced access pattern?

\begin{Verbatim}[xleftmargin=2em]
__global__ void histo_kernel(unsigned char *buf, long size,
                             unsigned int *histo) {
    int i = threadIdx.x + blockIdx.x * blockDim.x;
    int stride = ????
    while (i < size) {
        int pos = buf[i] - 'a';
        if (pos >= 0 && pos < 26)
            atomicAdd(&histo[pos/4], 1);
        i += stride;
    }
}
\end{Verbatim}

\vspace{4pt}
\texttt{int stride = }\,\widebox

\vspace{8pt}
% ---------------- Q3 ----------------
\question \textbf{[4 points]} A teammate proposes that in the kernel above, having every block
maintain a \texttt{\_\_shared\_\_} private histogram and merging into the global histogram at
the end will significantly improve performance. Briefly explain \emph{why} this typically
helps.\\[10pt]
\Ans \paragraphbox{2.6cm}

\vspace{8pt}
% ---------------- Q4 ----------------
\question \textbf{[13 points] Unified Memory}\\
You write the following kernel and host program on a \textbf{Pascal-class} GPU (CUDA 8+):

\begin{Verbatim}[xleftmargin=1em]
__global__ void mykernel(char *data) {
    data[1] = 'g';
}

void foo() {
    char *data;
    cudaMallocManaged(&data, 2);
    mykernel<<<1,1>>>(data);
    // no synchronize here
    data[0] = 'c';
    cudaFree(data);
}
\end{Verbatim}

\begin{parts}
  \part {}[3 points] On Pascal+, what happens when the CPU executes \texttt{data[0] = 'c'} while
        the kernel may still be running?\\[8pt]
        \Ans \paragraphbox{1.8cm}
  \part {}[3 points] What would have happened on a pre-Pascal device (e.g., Kepler) running the
        same code?\\[8pt]
        \Ans \paragraphbox{1.8cm}
  \part {}[3 points] You then increase the allocation to 64 GB on a 16 GB Pascal GPU. Will the
        \texttt{cudaMallocManaged} call succeed? Briefly justify.\\[8pt]
        \Ans \paragraphbox{1.8cm}
  \part {}[4 points] Name the API used to give the runtime a hint that data will be
        \textbf{read-mostly}, and describe what the runtime does as a result.\\[8pt]
        \Ans \paragraphbox{2.4cm}
\end{parts}

\newpage
% ---------------- Q5 ----------------
\question \textbf{[12 points] Tiled SGEMM Analysis}\\
Consider a square matrix multiply with \texttt{Width = 1024}.

\begin{parts}
  \part {}[3 points] For \texttt{TILE\_WIDTH = 16}, how many global memory loads does each
        \emph{thread block} perform per phase?\\[8pt]
        \Ans \smallbox
  \part {}[3 points] For \texttt{TILE\_WIDTH = 16}, how many phases (outer-loop iterations) does
        each block execute?\\[8pt]
        \Ans \smallbox
  \part {}[3 points] For \texttt{TILE\_WIDTH = 32}, how many bytes of shared memory does each
        block consume?\\[8pt]
        \Ans \smallbox
  \part {}[3 points] On an SM with 48 KB of shared memory, how many \texttt{TILE\_WIDTH = 32}
        blocks can be co-resident based on the shared-memory limit alone?\\[8pt]
        \Ans \smallbox
\end{parts}

\newpage
% ---------------- MC ----------------
\textbf{Multiple Choice. \textnormal{Please completely fill in the circle} $\bigcirc$
\textnormal{for your answer.} [12 points total]}

\vspace{6pt}
% Q6
\question {}[3 points] Which CUDA API call \emph{actually moves data} between processors in a
Unified Memory program (i.e., does not merely hint at future behavior)?
\begin{itemize}[leftmargin=2.5em,itemsep=4pt,label={}]
  \item \mc \texttt{cudaMemAdvise(..., cudaMemAdviseSetReadMostly, ...)}
  \item \mc \texttt{cudaMemAdvise(..., cudaMemAdviseSetPreferredLocation, ...)}
  \item \mc \texttt{cudaMemAdvise(..., cudaMemAdviseSetAccessedBy, ...)}
  \item \mc \texttt{cudaMemPrefetchAsync(...)}
\end{itemize}

% Q7
\question {}[3 points] An access pattern \texttt{A[expr]} in a CUDA kernel is best classified as
\textbf{coalesced} when:
\begin{itemize}[leftmargin=2.5em,itemsep=4pt,label={}]
  \item \mc \texttt{expr} is a multiple of \texttt{threadIdx.x}
  \item \mc \texttt{expr} = (something independent of \texttt{threadIdx.x}) + \texttt{threadIdx.x}
  \item \mc \texttt{expr} = \texttt{threadIdx.y * Width + threadIdx.x * Width}
  \item \mc \texttt{expr} = a function of \texttt{blockIdx.x} only
\end{itemize}

% Q8
\question {}[3 points] On a Kepler GPU using Hyper-Q with \textbf{32 hardware work queues}, two
processes share the GPU via MPS, and each process issues 5 streams. What is the most likely
performance effect for this configuration?
\begin{itemize}[leftmargin=2.5em,itemsep=4pt,label={}]
  \item \mc Full 32-way concurrency, no false dependencies
  \item \mc False dependencies/serialization due to more streams per client than channels
  \item \mc MPS automatically rebalances work across GPUs to avoid contention
  \item \mc Kernels from different processes always serialize, regardless of MPS
\end{itemize}

% Q9
\question {}[3 points] In a Dynamic Parallelism kernel, if every thread in a 256-thread block
executes the line \texttt{child<<<1,1>>>(...)} without any guard, how many child kernel
invocations are spawned by that block?
\begin{itemize}[leftmargin=2.5em,itemsep=4pt,label={}]
  \item \mc 1
  \item \mc 32 (one per warp)
  \item \mc 256
  \item \mc 0 — only the host can launch kernels
\end{itemize}

\newpage
% ---------------- T/F ----------------
\textbf{True or False. \textnormal{For the following questions, fill the answer box with (T)
True or (F) False.} [12 points total]}

\vspace{10pt}
\renewcommand{\arraystretch}{1.3}
\noindent
\begin{tabular}{@{}p{0.40\linewidth}c@{\hspace{0.5cm}}p{0.40\linewidth}c@{}}
10.~\texttt{atomicAdd} returns the new value (after the add). & \tfbox
 & 13.~\texttt{cudaDeviceSynchronize()} called inside a device kernel implies
   \texttt{\_\_syncthreads()} as well. & \tfbox \\[1.6cm]
11.~Memory coalescing rules apply to shared-memory accesses the same way they do to global
memory. & \tfbox
 & 14.~On a dual-IOH server, peer-to-peer (P2P) copy between two GPUs attached to different IOHs
   can still complete correctly. & \tfbox \\[1.6cm]
12.~In the tiled SGEMM kernel, the two \texttt{\_\_syncthreads()} calls per phase can be safely
removed for large enough block sizes. & \tfbox
 & 15.~\texttt{cudaMemAdviseSetAccessedBy} causes the data to immediately move to the indicated
   device. & \tfbox \\[1.2cm]
\end{tabular}

\vspace{16pt}
\setcounter{question}{15}
% ---------------- Q16 ----------------
\question {}[6 points] You write a tiled matrix-multiply with \texttt{TILE\_WIDTH = 16} for a
\textbf{non-square} input: $M$ is $j \times k$, $N$ is $k \times l$, $P$ is $j \times l$, with
$j = 1000$, $k = 800$, $l = 1500$.
\begin{parts}
  \part {}[2 points] What are the x and y dimensions of your grid?\\[8pt]
        \Ans \smallbox
  \part {}[2 points] Write the boundary condition a thread must check before loading its element
        of the \textbf{M tile} (use variables \texttt{Row, p, TILE\_WIDTH, tx, j, k}).\\[8pt]
        \Ans \widebox[12cm]
  \part {}[2 points] Write the boundary condition a thread must check before \textbf{writing}
        its \texttt{P} element (use \texttt{Row, Col, j, l}).\\[8pt]
        \Ans \widebox[12cm]
\end{parts}

\newpage
% ---------------- Q17 ----------------
\question {}[12 points] \textbf{Dynamic Parallelism trace.}\\
Consider the following device-side code:

\begin{Verbatim}[xleftmargin=1em]
__global__ void parent(int *out) {
    int tid = threadIdx.x;             // PC = A
    if (tid < 4) {                     // block of 32 threads
        child<<<1, 8>>>(out, tid);     // PC = B
    }
    cudaDeviceSynchronize();           // PC = C
    out[tid] += 1;                     // PC = D
}
\end{Verbatim}

\begin{parts}
  \part {}[3 points] How many \emph{child grids} are spawned by this block? Briefly justify.\\[8pt]
        \Ans \paragraphbox{2.0cm}
  \part {}[3 points] At PC = C, exactly which launches have finished from the perspective of
        thread \texttt{tid = 5}, which never executed the launch line?\\[8pt]
        \Ans \paragraphbox{2.0cm}
  \part {}[3 points] At PC = D, can thread \texttt{tid = 5} safely read \texttt{out[k]} written
        by a child launched by thread 0? Why or why not?\\[8pt]
        \Ans \paragraphbox{2.0cm}
  \part {}[3 points] If you replaced \texttt{cudaDeviceSynchronize()} with \texttt{\_\_syncthreads()},
        would the program still be correct? Briefly justify.\\[8pt]
        \Ans \paragraphbox{2.0cm}
\end{parts}

\vspace{14pt}
% ---------------- Q18 ----------------
\question {}[4 points] In a histogram kernel, two threads in the \emph{same} block simultaneously
attempt \texttt{atomicAdd(\&histo\_private[3], 1)} where \texttt{histo\_private} lives in shared
memory and currently holds 7. After both atomics complete, what is the value of
\texttt{histo\_private[3]}, and what values did the two threads receive as return values?\\[10pt]
\Ans \paragraphbox{2.6cm}

\end{questions}

\end{document}
